% Updated RQ6 Section Text
% Replace the existing RQ6 section with this updated version

\subsection{RQ6: Overhead}

We measure the overhead of our unsoundness detection framework across \emph{CBR sampling budgets} (5, 10, 20, 50)—i.e., the number of concrete inputs sampled by CBR per network—and model sizes (Small, Medium, Large).
CBR overhead is expected to scale linearly with budget, while BBL overhead depends on model size but remains constant per budget.
We report absolute overhead times (in milliseconds) and analyze the cost breakdown between CBR and BBL.

Table~\ref{tab:rq6-overhead-full} reports overhead across all configurations.
CBR time scales approximately linearly with sampling budget:
at Small model size, overhead increases from 0.50\,ms (budget=5) to 2.81\,ms (budget=50), confirming that sampling cost dominates CBR.
BBL time is independent of sampling budget and grows modestly with model size (0.09\,ms for Small, 0.13\,ms for Medium/Large), remaining significantly smaller than CBR across all configurations.

At the default configuration (budget=20), combined overhead (CBR+BBL) ranges from 1.28\,ms (Medium) to 1.57\,ms (Large).
For Small models, overhead at budget=50 (2.90\,ms) is approximately 5× that at budget=5 (0.59\,ms), consistent with linear CBR scaling.
The cost breakdown shows CBR accounts for 85--95\% of total overhead, with BBL contributing a small constant cost.

\finding{Our unsoundness detection framework exhibits predictable overhead scaling: CBR cost grows linearly with sampling budget (0.5--3.6\,ms for budgets 5--50), while BBL adds a small constant cost (0.09--0.13\,ms) independent of budget. At default settings (budget=20), total overhead is 1.3--1.6\,ms across model sizes.}
